\documentclass[paper=screen,mode=present,style=zysimple]{powerdot}

\usepackage{fancyvrb,verbatim,listings}
\usepackage{proof}
\usepackage[all]{xypic,xy}

\usepackage{latexsym,stmaryrd,amsmath,amsfonts}

%% \newtheorem{definition}{Definition}
%% \newtheorem{theorem}{Theorem}
%% \newtheorem{lemma}{Lemma}
\newtheorem{prop}{Proposition}
%% \newtheorem{corollary}{Corollary}
%% \newenvironment{proof}{{\sc Proof . }}{{\hfill $\square$}}
%% \newcommand{\myqed}{\hfill $\square$}
%% \newcommand{\ped}{\vspace{-1.6em}}
%% \newcommand{\note}[1]{\noindent\fbox{\parbox{\textwidth}{\em #1}}}

\newcommand{\abind}[2]{\mathtt{bind} \ (\, {#1} \,)  \ \mathtt{in} \ {#2}}
\newcommand{\acase}{\mathtt{case}}
\newcommand{\adv}{\mathcal{A}}
\newcommand{\aspecto}[1]{\xrightarrow{\; {#1} \;}}
\newcommand{\asplinear}{\mathfrak{l}}
\newcommand{\aspmodal}{\mathfrak{m}}
\newcommand{\aspnormal}{\mathfrak{n}}
%
\newcommand{\betareduce}{\leadsto_\beta}
\newcommand{\betareduceto}{\twoheadrightarrow_\beta}
\newcommand{\bit}{\mathtt{B}}
\newcommand{\bitone}{\mathtt{B}_1}
\newcommand{\bitzero}{\mathtt{B}_0}
%
\newcommand{\caseto}{\triangleright}
\newcommand{\cbind}[3]{\mathtt{bind} \ {#1} = {#2} \ \mathtt{in} \ {#3}}
\newcommand{\ccase}[4]{\mathtt{case} \; {#1} \; \mathtt{of} \; \cnil \caseto {#2}; \;
                                          \bitzero x_0 \caseto {#3}; \; \bitone x_1 \caseto {#4}}
\newcommand{\cequiv}{\simeq}
\newcommand{\cif}[3]{\mathtt{if} \; {#1} \; \mathtt{then} \; {#2} \; \mathtt{else} \; {#3}}
\newcommand{\ctrue}{\mathtt{true}}
\newcommand{\cfalse}{\mathtt{false}}
\newcommand{\cfix}{\mathtt{fix}\,}
\newcommand{\cfixin}[2]{\mathtt{fix}\; {#1}\; \mathtt{in} \; {#2}}
\newcommand{\cfold}{\mathtt{fold}\,}
\newcommand{\chd}{\mathtt{hd}}
\newcommand{\cinj}{\mathtt{inj}}
\newcommand{\clift}[1]{\lfloor{#1}\rfloor}
\newcommand{\clet}[3]{\mathtt{let} \ {#1} = {#2} \ \mathtt{in} \ {#3}}
\newcommand{\cletst}[3]{\mathtt{liftlet} \ {#1} = {#2} \ \mathtt{in} \ {#3}}
\newcommand{\cletrec}[3]{\mathtt{letrec} \ {#1} = {#2} \ \mathtt{in} \ {#3}}
\newcommand{\cnil}{\mathtt{nil}}
\newcommand{\conc}{\! \bullet \!}
\newcommand{\concat}{\pmb{\mathit{conc}}}
\newcommand{\context}{\mathbb{C}}
\newcommand{\cop}{\;\mathtt{op}\;}
\newcommand{\cpair}[1]{\langle{#1}\rangle}
\newcommand{\cproj}{\mathtt{proj}}
\newcommand{\crand}{\mathtt{rand}}
\newcommand{\crec}{\mathtt{rec}}
\newcommand{\croll}{\mathtt{roll}\,}
\newcommand{\crule}[1]{\mbox{{\textsl{#1}}}}
\newcommand{\cslrtype}{\Vdash}
\newcommand{\csucc}{\mathtt{S}}
\newcommand{\ctl}{\mathtt{tl}}
\newcommand{\cut}{\pmb{\mathit{cut}}}
\newcommand{\cval}{\mathtt{return}}
\newcommand{\cwhile}[2]{\mathtt{while}\; {#1} \; \mathtt{do} \; {#2} \; \mathtt{od}}
\newcommand{\cvalof}{\pmb{valueof}}
\newcommand{\cunfold}{\mathtt{unfold}\,}
\newcommand{\cunit}{\mathtt{unit}}
\newcommand{\cunroll}{\mathtt{unroll}\,}
\newcommand{\cupdate}{\pmb{update}}
\newcommand{\czero}{\mathtt{O}}
%
\newcommand{\deduces}{\vdash}
\newcommand{\defeq}{\stackrel{\mathrm{def}}{=}}
\newcommand{\delay}[1]{\left\lfloor{#1}\right\rfloor}
\newcommand{\distr}{\mathcal{D}}
%
\newcommand{\eagerreduce}{\xrightarrow{\mathrm{eager}}}
\newcommand{\Eagerreduce}{\stackrel{\mathrm{eager}}{\Reduce}}
\newcommand{\eqtest}{\stackrel{\mbox{\tiny ?}}{=}}
\newcommand{\Eval}[1]{{\left\llbracket{#1}\right\rrbracket}}
\newcommand{\evalcontext}{\mathbb{E}}
\newcommand{\Example}{\emgreen{\pmb{\textsc{Example}\;}}}
\newcommand{\expansion}[1]{\tilde{{#1}}}
%
\newcommand{\ford}{\mathrm{ord}}
%
\newcommand{\getval}{\leftarrow}
%
\newcommand{\head}{\pmb{\mathit{hd}}}
\newcommand{\termhole}{[\;]}
\newcommand{\hoaretriple}[3]{\{{#1}\} \; {#2} \; \{{#3}\}}
%
\newcommand{\imply}{\Rightarrow}
\newcommand{\indist}{\simeq}
\newcommand{\iteration}{\pmb{\mathit{it}}}
%
\newcommand{\lazyreduce}{\xrightarrow{\mathrm{lazy}}}
\newcommand{\Lazyreduce}{\lazyreduce\!\!\!\!\!\to}
\newcommand{\ldot}{\, . \, }
\newcommand{\linearto}{\multimap}
\newcommand{\logrel}{(\mathcal{R}_\tau)_{\mbox{$\tau$ type}}}
\newcommand{\lr}[1]{\mathcal{R}_{#1}}
\newcommand{\lub}{\sqcup}
%
\newcommand{\metalambda}{\underline{\lambda}}
\newcommand{\modal}{\square}
%
\newcommand{\nat}{\mathbb{N}}
\newcommand{\nReduce}{\nrightarrow\!\!\!\!\to}
\newcommand{\nreduceto}{\nrightsquigarrow}
\newcommand{\nullenv}{\varnothing}
%
\newcommand{\orc}{\mathcal{O}}
\newcommand{\oslrtype}{\Vdash}
%
\newcommand{\porder}{\sqsubseteq}
\newcommand{\prefix}{\pmb{\mathit{pref}}}
\newcommand{\proba}{\mathbf{Pr}}
%
\newcommand{\randlen}{\pmb{\mathit{rl}}}
\newcommand{\randnum}{\pmb{\mathit{rn}}}
\newcommand{\randstr}{\pmb{\mathit{rs}}}
\newcommand{\reduce}{\leadsto}
\newcommand{\reduceto}{\rightsquigarrow\hspace*{-6px}\shortrightarrow}
\newcommand{\Reduce}{\to\!\!\!\to}
\newcommand{\Reduceto}{\Downarrow}
%
\newcommand{\set}[1]{\{{#1}\}}
\newcommand{\setbits}{\mathbb{B}}
\newcommand{\slrtype}{\vdash}
\newcommand{\splt}{\pmb{\mathit{split}}}
\newcommand{\suffix}{\pmb{\mathit{suff}}}
\newcommand{\sysname}{\pmb{\mathcal{C}}}
%
\newcommand{\tail}{\pmb{\mathit{tl}}}
\newcommand{\tbits}{\mathsf{Bits}}
\newcommand{\tbool}{\mathsf{Bool}}
\newcommand{\tnat}{\mathsf{Nat}}
\newcommand{\tcomp}{\mathsf{T}}
\newcommand{\tlift}[1]{\lfloor{#1}\rfloor}
\newcommand{\tlist}{\mathsf{List}}
\newcommand{\tself}{\mathsf{self}}
\newcommand{\tstore}{\mathsf{Store}}
\newcommand{\tstream}{\mathsf{Stream}}
\newcommand{\tloc}{\mathsf{Loc}}
\newcommand{\ttcase}{\mathtt{case}}
\newcommand{\ttif}{\mathtt{if}}
\newcommand{\ttthen}{\mathtt{then}}
\newcommand{\ttelse}{\mathtt{else}}
\newcommand{\ttin}{\mathtt{in}}
\newcommand{\ttlet}{\mathtt{let}}
\newcommand{\ttof}{\mathtt{of}}
\newcommand{\tunit}{\mathsf{Unit}}
\newcommand{\Type}{\mathcal{T}\!\mathit{ype}}
%
\newcommand{\vbool}{\pmb{\mathbb{B}}}
\newcommand{\vff}{\mathfrak{f}\!\mathfrak{f}}
\newcommand{\vnat}{\pmb{\mathrm{N}}}
\newcommand{\vnil}{\epsilon}
\newcommand{\vtt}{\mathfrak{t}\!\mathfrak{t}}
%
\newcommand{\weakpre}{\pmb{wp}}




\pdsetup{lf={Theory and Practice of Functional Programming}}

\title{{\small \textsc{Theory and Practice of Functional Programming}}
\\[12pt]
{\Large \color{TitleColor} Simply Typed Lambda-Calculi (II)}
}
\author{\small \textbf{Dr. ZHANG Yu}
  \\[2pt] \small Institute of Software, Chinese Academy of Sciences 
  \\[5pt] \includegraphics[height=.04\slideheight]{iscas.eps} 
  \qquad \qquad 
  \includegraphics[height=.04\slideheight]{cas.eps}
}
\date{GUCAS, Beijing}

\newcommand{\RED}{\mathsf{RED}}

\begin{document}

\maketitle 

\begin{slide}{Introduction}
\begin{itemize}
\item PCF --- {\em Programming Computable Functions}, mainly designed for program analysis 
  and reasoning, rather than practical use.
\begin{itemize}
\item Syntax and type system 
\item Reduction and operational semantics
\item Expressive power
\end{itemize}
\end{itemize}
\end{slide}

\begin{slide}{Base types}
\begin{itemize}
\item Unit type: $\tunit$, which contains a single element $*$.
\item Natural numbers: 
\begin{itemize}
\item Constants (values): $0, 1, 2, \ldots$.
\item Arithmetic operations: $e_1 + e_2, e_1 * e_2, \ldots$.
\end{itemize}
$$\small
\begin{prooftree}
\Gamma \vdash e_1 : \tnat
\quad \Gamma \vdash e_2 : \tnat
\quad \mathtt{op} \in \set{+, -, *, /, \ldots}
\justifies 
\Gamma \vdash e_1 \;\mathtt{op}\; e_2 : \tnat 
\end{prooftree}
\vspace*{-2em}
$$
\item Booleans:
\begin{itemize}
\item Boolean constants (values): $\ctrue, \cfalse$.
\item Equality test (of natural numbers): $e_1 \eqtest e_2$.
\item Conditional: $\cif{e}{e_1}{e_2}$. 
\end{itemize}
$$\small
\begin{prooftree}
\Gamma \vdash e_1 : \tnat
\quad \Gamma \vdash e_2 : \tnat
\justifies 
\Gamma \vdash e_1 \eqtest e_2 : \tbool
\end{prooftree}
\qquad 
\begin{prooftree}
\Gamma \vdash e : \tbool 
\quad \Gamma \vdash e_1 : \tau 
\quad \Gamma \vdash e_2 : \tau 
\justifies 
\Gamma \vdash \cif{e}{e_1}{e_2} : \tau 
\end{prooftree}
$$
\end{itemize}
\end{slide}

\begin{slide}{More types}
\begin{itemize}
\item Product types: $\tau_1 \times \tau_2$ 
% \begin{itemize}
% \item Pairs: $\cpair{e_1, e_2}$.
% \item Projections: $\cproj_i \cpair{e_1, e_2} = e_i$ ($i = 1, 2$).
% \end{itemize}
\[\small
\begin{prooftree}
\Gamma \vdash e_1 : \tau_1 
\quad 
\Gamma \vdash e_2 : \tau_2 
\justifies 
\Gamma \vdash \cpair{e_1, e_2} : \tau_1 \times \tau_2
\end{prooftree}
\qquad 
\begin{prooftree}
\Gamma \vdash e : \tau_1 \times \tau_2
\justifies 
\Gamma \vdash \cproj_i\ e : \tau_i
\end{prooftree}
\vspace*{-1.2em}
\]
\item Sum types: $\tau_1 + \tau_2$
\begin{itemize}
\item A value of $\tau_1 + \tau_2$ can be a value of either $\tau_1$ or $\tau_2$, distinguished by ``tags'' ($\cinj_1, \cinj_2$).
\item We use $\mathtt{case}$ to retrieve the value (removing tags).
\end{itemize}
\[\small
\begin{prooftree}
\Gamma \vdash e : \tau_i
\justifies
\Gamma \vdash \cinj_i \ e : \tau_1 + \tau_2
\end{prooftree}
\qquad 
\begin{prooftree}
\Gamma, x_i : \tau_1 \vdash e_i : \tau (i = 1, 2)
\quad 
\Gamma \vdash e : \tau_1 + \tau_2
\justifies
\Gamma \vdash 
\mathtt{case} \; e \; \mathtt{of} \; \cinj_1 x_1 \; \mathtt{in} \; e_1 ; \cinj_2 x_2 \; \mathtt{in} \; e_2 : \tau
\end{prooftree}
% \vspace*{-0.5em}
\] 
Boolean type can be defined using $\tunit$ and sum: $\tbool \defeq \tunit + \tunit$. 
\vspace*{-0.5em}
\[
\begin{array}{c}
\ctrue \defeq \cinj_1 \ * ,
\qquad
\cfalse \defeq \cinj_2 \ *
\\[0.2em]
\cif{e}{e_1}{e_2} \defeq 
\mathtt{case} \; e \; \mathtt{of} \; \cinj_1 \_ \; \mathtt{in} \; e_1 ; \cinj_2 \_ \; \mathtt{in} \; e_2
\end{array}
\]
\end{itemize}
\end{slide}

\begin{slide}[method=direct]{Functions and recursions}
%% \begin{verbatim}
%% int fact (x : int) {
%%   int y;
%%   if (x == 0) then y = 1 else y = x * fact(x-1);
%%   return y;
%% }
%% \end{verbatim}
\begin{itemize}
\item Function type $\tau \to \tau'$
$$ \small
\begin{prooftree}
\Gamma, x: \tau  \vdash e : \tau'
\justifies
\Gamma \vdash \lambda x : \tau \ldot e : \tau \to \tau'
\end{prooftree}
\qquad 
\begin{prooftree}
\Gamma \vdash e_1 : \tau \to \tau'
\quad 
\Gamma \vdash e_2 : \tau
\justifies
\Gamma \vdash e_1 e_2 : \tau'
\end{prooftree}
$$
\item PCF introduces explicitly an operator for fix-point: $\cfix_\tau$.
\begin{itemize}
\item If $F$ is a function of type $\tau \to \tau$, then the fix-point of $F$ 
  is the value $x : \tau$ such that \emred{$F(x) = x$}, written as \emblue{$\cfix_\tau F$}. 
\item A recursive definition $f : \tau = e$ (with $f$ occurs freely in $e$) indeed defines a fix-point of 
  \emred{$\lambda f : \tau \ldot e$}.
\item Example: factorization can be defined recursively: \vspace*{-0.5em}
$$f = \lambda x : \tnat \ldot \cif{x \eqtest 0}{1}{x * f(x-1)} $$
With fix-point operator in PCF, we shall write the definition as \vspace*{-0.5em}
$$fact \defeq \cfix_{\tnat \to \tnat} (\lambda f \ldot \lambda x \ldot \cif{x \eqtest 0}{1}{x * f(x-1)})$$
\vspace*{-1.5em}
\end{itemize}
\end{itemize}
\end{slide}

\begin{slide}{Syntax of PCF}
\begin{itemize}
\item Types:
\[
\tau_1, \tau_2, \ldots ::= \tnat \mid \tbool \mid \tau_1 \times \tau_2 \mid \tau_1 + \tau_2 \mid \tau_1 \to \tau_2
\vspace*{-2em}
\]
% where $b$ ranges over a set of base types $\{ \tbool, \tnat, \ldots \}$.
%% \item Order of types: 
%% \[
%% \begin{array}{l}
%% \ford(b) = 0 , \qquad \mbox{if $b$ is a base type} \\
%% \ford(\tau_1 \times \tau_2) = \max(\ford(\tau_1), \ford(\tau_2)) \\
%% \ford(\tau_1 + \tau_2) = \max(\ford(\tau_1), \ford(\tau_2)) \\
%% \ford(\tau_1 \to \tau_2) = \max(\ford(\tau_1)+1, \ford(\tau_2))
%% \end{array}
%% \]
\item Expressions (terms): \vspace*{-0.5em}
\[\small
\begin{array}{lrll}
e_1, e_2, \ldots & ::= &
x & \mbox{variables} 
\\ & \mid & 
0, 1, 2,  \ldots & \mbox{natural numbers}
\\ & \mid & 
e_1 \; \mathtt{op} \; e_2 & \mbox{arithmetic operations} 
\\ & \mid & 
\ctrue \mid \cfalse & \mbox{booleans}
\\ & \mid & 
e_1 \eqtest e_2 & \mbox{equality test} 
\\ & \mid & 
\cif{e}{e_1}{e_2} & \mbox{conditional}
\\ & \mid & 
\cpair{e_1, e_2} & \mbox{pairing}
\\ & \mid & 
\cproj_1 e \mid \cproj_2 e & \mbox{projections}
\\ & \mid & 
\cinj_1 e \mid \cinj_2 e & \mbox{injections}
\\ & \mid & 
\mathtt{case} \; e \; \mathtt{of} \; \cinj_1 x_1 \; \mathtt{in} \; e_1 ; 
& \mbox{case distinction}
\\ & & \hspace*{4.5em}
\cinj_2 x_2 \; \mathtt{in} \; e_2 & 
\\ & \mid & 
\lambda x : \tau \ldot e \mid e_1 \, e_2 & \mbox{functions and applications}
\\ & \mid & 
\cfix_\tau & \mbox{recursion (fix point)}
\end{array}
\]
\end{itemize}
\end{slide}

\begin{slide}{PCF typing rules}
\vspace*{-1em}
\[\small
\begin{array}{c}
\begin{prooftree}
x : \tau \in \Gamma
\justifies
\Gamma \vdash x : \tau 
\end{prooftree}
\qquad
\begin{prooftree}
n = 0, 1, 2, \ldots 
\justifies
\Gamma \vdash n : \tnat 
\end{prooftree}
\qquad 
\begin{prooftree}
\Gamma \vdash e_1 : \tnat 
\quad \Gamma \vdash e_2 : \tnat 
\quad \mathtt{op} \in \set{+, -, * , /}
\justifies
\Gamma \vdash e_1 \;\mathtt{op}\; e_2 : \tnat 
\end{prooftree}
\\[2em]
\begin{prooftree}
\justifies
\Gamma \vdash \ctrue : \tbool
\end{prooftree}
\qquad 
\begin{prooftree}
\justifies
\Gamma \vdash \cfalse : \tbool
\end{prooftree}
\qquad 
\begin{prooftree}
\Gamma \vdash e_1 : \tnat 
\qquad 
\Gamma \vdash e_2 : \tnat 
\justifies
\Gamma \vdash e_1 \eqtest e_2 : \tbool
\end{prooftree}
\\[2em]
\begin{prooftree}
\Gamma \vdash e : \tbool
\qquad 
\Gamma \vdash e_1 : \tau
\qquad 
\Gamma \vdash e_2 : \tau 
\justifies
\Gamma \vdash \cif{e}{e_1}{e_2} : \tau 
\end{prooftree}
\qquad 
\begin{prooftree}
\Gamma \vdash e_1 : \tau_1
\qquad 
\Gamma \vdash e_2 : \tau_2 
\justifies
\Gamma \vdash \cpair{e_1, e_2} : \tau_1 \times \tau_2
\end{prooftree}
\\[2em]
\begin{prooftree}
\Gamma \vdash e : \tau_1 \times \tau_2
\justifies
\Gamma \vdash \cproj_1 e : \tau_1
\end{prooftree}
\qquad
\begin{prooftree}
\Gamma \vdash e : \tau_1 \times \tau_2
\justifies
\Gamma \vdash \cproj_2 e : \tau_2 
\end{prooftree}
\qquad
\begin{prooftree}
\Gamma \vdash e : \tau_1
\justifies
\Gamma \vdash \cinj_1 \ e : \tau_1 + \tau_2
\end{prooftree}
\qquad 
\begin{prooftree}
\Gamma \vdash e : \tau_2
\justifies
\Gamma \vdash \cinj_2 \ e : \tau_1 + \tau_2
\end{prooftree}
\\[2em]
\begin{prooftree}
\Gamma, x_1 : \tau_1 \vdash e_1 : \tau 
\qquad 
\Gamma, x_2 : \tau_2 \vdash e_2 : \tau 
\qquad 
\Gamma \vdash e : \tau_1 + \tau_2
\justifies
\Gamma \vdash 
\mathtt{case} \; e \; \mathtt{of} \; \cinj_1 x_1 \; \mathtt{in} \; e_1 ; \cinj_2 x_2 \; \mathtt{in} \; e_2 : \tau
\end{prooftree}
\\[2em]
\begin{prooftree}
\Gamma, x: \tau  \vdash e : \tau'
\justifies
\Gamma \vdash \lambda x : \tau \ldot e : \tau \to \tau'
\end{prooftree}
\qquad 
\begin{prooftree}
\Gamma \vdash e_1 : \tau \to \tau'
\qquad 
\Gamma \vdash e_2 : \tau
\justifies
\Gamma \vdash e_1 e_2 : \tau'
\end{prooftree}
\qquad
\begin{prooftree}
\Gamma \vdash e : \tau \to \tau
\justifies
\Gamma \vdash \cfix_\tau e : \tau
\end{prooftree}
\end{array}
\]
\end{slide}

\begin{slide}{Type checking}
\begin{itemize}
\item A procedure of checking the validity of a type assertion, by typing rules and 
  induction on expression structure.
\item Example: 
  $\cfix_{\tnat \to \tnat} \lambda f : \tnat \to \tnat \ldot \lambda x : \tnat \ldot 
\cif{x \eqtest 0}{1}{x * f(x-1)}$
\end{itemize}
\vspace*{-0.5em}
\tiny
\[
\begin{prooftree}
\begin{prooftree}
\begin{prooftree}
\begin{prooftree}
\begin{prooftree}
\begin{prooftree}
\justifies 
\Gamma\vdash  x : \tnat
\end{prooftree}
\quad 
\begin{prooftree}
\justifies 
\Gamma \vdash  0 : \tnat
\end{prooftree}
\justifies 
\Gamma \vdash  x \eqtest 0 : \tbool
\end{prooftree}
\quad
\begin{prooftree}
\justifies 
\Gamma \vdash 1 : \tnat 
\end{prooftree}
\quad 
\begin{prooftree}
\begin{prooftree}
\justifies 
\Gamma \vdash  x : \tnat
\end{prooftree}
\quad 
\begin{prooftree}
\begin{prooftree}
\justifies 
\Gamma \vdash  f : \tnat \to \tnat
\end{prooftree}
\quad 
\begin{prooftree}
\begin{prooftree}
\justifies 
\Gamma \vdash  x : \tnat
\end{prooftree}
\quad 
\begin{prooftree}
\justifies 
\Gamma \vdash  1 : \tnat
\end{prooftree}
\justifies 
\Gamma \vdash  x-1 : \tnat
\end{prooftree}
\justifies 
\Gamma \vdash  f(x-1) : \tnat
\end{prooftree}
\justifies 
\Gamma \vdash  x * f(x-1) : \tnat 
\end{prooftree}
\justifies 
f : \tnat \to \tnat, x : \tnat \vdash  
\cif{x \eqtest 0}{1}{x * f(x-1)} : \tnat
\end{prooftree}
\justifies 
f : \tnat \to \tnat \vdash \lambda x : \tnat \ldot 
\cif{x \eqtest 0}{1}{x * f(x-1)} : \tnat \to \tnat
\end{prooftree}
\justifies 
\vdash \lambda f : \tnat \to \tnat \ldot \lambda x : \tnat \ldot 
\cif{x \eqtest 0}{1}{x * f(x-1)} : (\tnat \to \tnat) \to (\tnat \to \tnat)
\end{prooftree}
\justifies
\vdash \cfix_{\tnat \to \tnat} \lambda f : \tnat \to \tnat \ldot \lambda x : \tnat \ldot 
\cif{x \eqtest 0}{1}{x * f(x-1)} : \tnat \to \tnat
\end{prooftree}
\]
\end{slide}

\begin{slide}{PCF reduction}
\begin{itemize}
\item $\beta$-reduction: $(\lambda x : \tau \ldot e) e' \reduce e [e' / x]$.
\item Other reductions:
\[
\begin{array}{l}
2 + 3 \reduce 5, 
\quad 12 \eqtest 23 \reduce \cfalse ,
\quad \ldots
\\[0.2em]
\cif{\ctrue}{e_1}{e_2} \reduce e_1, 
\quad \cif{\cfalse}{e_1}{e_2} \reduce e_2
\\[0.2em]
\cproj_1 \cpair{e_1, e_2} \reduce e_1, 
\quad \cproj_2 \cpair{e_1, e_2} \reduce e_2
\\[0.2em]
\mathtt{case} \; \emblue{\cinj_1 e} \; \mathtt{of} \; \cinj_1 x_1 \; \mathtt{in} \; e_1 ; \cinj_2 x_2 \; \mathtt{in} \; e_2 
\reduce \emblue{e_1[e / x_1]}
\\[0.2em]
\mathtt{case} \; \emblue{\cinj_2 e} \; \mathtt{of} \; \cinj_1 x_1 \; \mathtt{in} \; e_1 ; \cinj_2 x_2 \; \mathtt{in} \; e_2 
\reduce \emblue{e_2[e / x_2]}
\end{array}
\vspace*{-1em}
\]
\item Reduction of fix-pointer: \vspace*{-0.5em}
$$\emred{\cfix_\tau \reduce \lambda f : \tau \to \tau \ldot f (\cfix_\tau f) } \vspace*{-0.5em}$$
PCF reductions may not terminate: \vspace*{-0.5em}
\[\small
\begin{array}{ll}
\cfix \, \lambda x \ldot x 
& \reduce (\lambda f \ldot \cfix \, f (\cfix \, f)) (\lambda x \ldot x) 
\\ & 
\reduce (\lambda x \ldot x) (\cfix \, \lambda x \ldot x)
\\ & 
\reduce \emblue{\cfix \, \lambda x \ldot x}
\reduce \cdots 
\end{array}
\]
\end{itemize}
\end{slide}

\begin{slide}{PCF reduction}
\begin{itemize}
\item Example:
\vspace*{-0.5em}
$$fact \defeq  \cfix_{\tnat \to \tnat} (\lambda f \ldot \lambda x \ldot \cif{x \eqtest 0}{1}{x * f(x-1)}) $$
$$\small
\begin{array}{ll}
& fact \; (n+1) 
\\ = & 
\cfix_{\tnat \to \tnat} (\lambda f \ldot \lambda x \ldot \cif{x \eqtest 0}{1}{x * f(x-1)}) (n+1)
\\ \reduce & 
(\lambda g \ldot g (\cfix_{\tnat \to \tnat} g)) (\lambda f \ldot \lambda x \ldot \cif{x \eqtest 0}{1}{x * f(x-1)}) (n+1)
\\ \reduce & 
Fact (\cfix_{\tnat \to \tnat} Fact) \; (n+1)
\\ & 
\mbox{(where $Fact = \lambda f \ldot \lambda x \ldot \cif{x \eqtest 0}{1}{x * f(x-1)}$)}
\\ = & 
(\lambda f \ldot \lambda x \ldot \cif{x \eqtest 0}{1}{x * f(x-1)}) fact \; (n+1)
\\ \reduce & 
(\lambda x \ldot \cif{x \eqtest 0}{1}{x * fact(x-1)}) (n+1)
\\ \reduce &
\cif{n+1 \eqtest 0}{1}{(n+1) * fact(n+1-1)} 
\\ \reduce &
\cif{\cfalse}{1}{(n+1) * fact(n+1-1)} 
\\ \reduce & 
(n+1) * fact(n+1-1) 
\end{array}
\vspace*{-1em}
$$
\item A syntactic abbreviation for recursion 
\vspace*{-0.8em}
\[
\begin{array}{l}
\emblue{\mathtt{letrec}} \; fact \emblue{=} 
\lambda x : \tnat \ldot \cif{x \eqtest 0}{1}{x * fact(x-1)} \\
\emblue{\mathtt{in}} \; fact \; (n+1)
\end{array}
\]
\end{itemize}
\end{slide}

\begin{slide}{PCF reduction}
\begin{itemize}
\item Multi-step reduction $\reduceto$: the reflexive (w.r.t. $\alpha$-equivalence) 
and transitive closure of $\reduce$. 
\vspace*{-0.5em}
\[
e \reduceto e' \mbox{ if } e =_\alpha e' \mbox{ or } e \reduce e'' \;\&\; e'' \Reduce e'
\vspace*{-2.5em}
\]
\item PCF reductions are \emblue{non-deterministic}. 
\vspace*{-0.5em}
\[
\begin{array}{l}
\cpair{12 + 47, 5 * 6} \reduce \cpair{59, 5 * 6} \reduce \cpair{59, 30}
\\
\cpair{12 + 47, 5 * 6} \reduce \cpair{12 + 47 , 30} \reduce \cpair{59, 30}
\end{array}
\vspace*{-2em}
\]
\item \empurple{PCF reduction system is confluent} --- it satisfies Church-Rosser property.
\begin{itemize}
\item No PCF term can have more than one normal form.
\item Inconsistent reduction system with \empurple{$\cfix e \reduce e (\cfix e)$} and $\eta$-reduction.
\[
\begin{array}{l}
\lambda x \ldot \cfix \, x \reduce_\eta \cfix  \\[0.6em]
\lambda x \ldot \cfix \, x \reduce_{\cfix} \lambda x \ldot x (\cfix \, x) 
\reduce_{\cfix} \lambda x \ldot x ( x (\cfix \, x)) \reduce \ldots
\end{array}
\]
\end{itemize}
\end{itemize}
\end{slide}

\begin{slide}{PCF values and canonical form}
\begin{itemize}
\item PCF values are \empurple{closed expressions that cannot be reduced any more}. 
\item PCF values are in \emblue{canonical forms}:
\begin{itemize}
\item Boolean values: $\ctrue$, $\cfalse$.
\item Natural number values: $0, 1, 2, \ldots$. 
\item Product values: $\cpair{e_1, e_2}$.
\item Injective values: $\cinj_1 e, \cinj_2 e$.
\item Function values: $\lambda x \ldot e$.
\end{itemize}
\item \empurple{Every closed normal PCF expression must be in the canonical form}. \\[0.3em]
{\em Proof.} By induction on the structure of PCF expressions.
\end{itemize}
\end{slide}

\begin{slide}{Reduction strategy}
\begin{itemize}
\item Restrict the reduction system so that reductions are \emred{deterministic} and consistent with 
  the original system. 
% \\ A reduction strategy defines a (partial) function from terms to terms.
% \\[0.3em]
\item Different reduction strategies
\[
\begin{array}{l@{\qquad}l}
(\lambda x: \tau \ldot e_1) e_2 \reduce e_1 [e_2 / x] 
& \mbox{Lazy evaluation, a.k.a, call-by-name} \\[0.8em]
(\lambda x: \tau \ldot e_1) e_2 \reduce (\lambda x: \tau \ldot e_1') e_2 
& \mbox{Function optimization} \\[0.8em]
(\lambda x: \tau \ldot e_1) e_2 \reduce (\lambda x: \tau \ldot e_1) e_2'
& \mbox{Eager evaluation, a.k.a, call-by-value}
\end{array}
\]
\end{itemize}
\end{slide}

\begin{slide}{A reduction system with lazy evaluation}
If the top-level term is a redex, it will be reduced first: if $e \reduce e'$, then $e \lazyreduce e'$.
\vspace*{-0.5em}
\[\small
\begin{array}{c}
\emred{
\begin{prooftree}
e_1 \lazyreduce e_1' 
\justifies
e_1 e_2  \lazyreduce e_1' e_2
\end{prooftree}
}
\qquad
\begin{prooftree}
e_1 \lazyreduce e_1'
\quad \mathtt{op} \in \set{+, -, *, /, \eqtest}
\justifies
e_1 \;\mathtt{op}\; e_2 \lazyreduce e_1' + e_2
\end{prooftree}
\\[2em] 
\begin{prooftree}
e \lazyreduce e' 
\quad n \in \set{0, 1, 2, \ldots}
\quad \mathtt{op} \in \set{+, -, *, /, \eqtest}
\justifies
n \;\mathtt{op}\; e \lazyreduce n + e'
\end{prooftree}
\\[2em]
\begin{prooftree}
e \lazyreduce e' 
\justifies
\cif{e}{e_1}{e_2} \lazyreduce \cif{e'}{e_1}{e_2}
\end{prooftree}
\qquad
\begin{prooftree}
e \lazyreduce e' 
\justifies
\cproj_i \, e \lazyreduce \cproj_i \, e'
\end{prooftree}
\\[2em]
\begin{prooftree}
e \lazyreduce e'
\justifies
\mathtt{case} \; e \; \mathtt{of} \; \cinj_1 x_1 \; \mathtt{in} \; e_1 ; \cinj_2 x_2 \; \mathtt{in} \; e_2 
\lazyreduce \mathtt{case} \; e' \; \mathtt{of} \ldots
\end{prooftree}
\end{array}
\vspace*{-1em}
\]
\begin{itemize}\small
\item Consistency with the PCF reduction system: if $e \Lazyreduce e'$, then $e \reduceto e'$
\item Incompleteness: lazy normal form is not necessarily PCF-normal form.
\end{itemize}
\end{slide}

%% \begin{slide}{Eager PCF}
%% \begin{itemize}
%% \item Also known as \emblue{call-by-value} reduction. \\
%% Values: $\ctrue, \cfalse, 0, 1, \ldots, \cpair{v_1, v_2}, \lambda x \ldot e, \delay{e}$ 
%% ($\delay{e} \defeq \lambda x : \tau \ldot e x$)
%% \\[0.4em]
%% \item Eager reductions ($v, v_1, v_2$ are values): \vspace*{-0.5em}
%% \[
%% \begin{array}{l}
%% 0 + 0 \eagerreduce 0, \; 0 + 1 \eagerreduce 1, \; \ldots \\[0.3em]
%% n \eqtest n \eagerreduce \ctrue, \; n \eqtest m \eagerreduce \cfalse (n \neq m) \\[0.3em]
%% \cif{\ctrue}{e_1}{e_2} \eagerreduce e_1 , \; 
%% \cif{\cfalse}{e_1}{e_2} \eagerreduce e_2 \\[0.3em]
%% \emblue{(\lambda x: \tau \ldot e) v \eagerreduce e [v / x]} \\[0.3em]
%% \cproj_i \cpair{v_1, v_2} \eagerreduce v_i \\[0.3em]
%% \emblue{\cfix v \eagerreduce v (\delay{\cfix v}) }
%% \end{array}
%% \vspace*{-1em}
%% \]
%% \item Evaluation context \vspace*{-0.6em}
%% \begin{eqnarray*}
%% \evalcontext_e & ::= & 
%% \termhole \mid \evalcontext_e + e \mid n + \evalcontext_e \mid \evalcontext_e \eqtest e 
%% \mid n \eqtest \evalcontext_e \mid \cif{\evalcontext_e}{e_1}{e_2} 
%% \\ & & 
%% \mid \cproj_i \, \evalcontext_e \mid \cpair{\evalcontext_e, e} \mid \cpair{v, \evalcontext_e}
%% \mid \evalcontext_e \, e \mid \emblue{v \, \evalcontext_e}
%% \end{eqnarray*}
%% \end{itemize}
%% \end{slide}

\begin{slide}{Lazy and eager evaluation}
\begin{itemize}
\item Lazy evaluation may lead to \emred{low efficient} computation, as it often produces multiple copies of 
  a expression, which will be executed (reduced) for multiple times. 
\vspace*{-0.5em}
\begin{itemize}
\item In practical FPL with lazy evaluation (e.g., Haskell), duplication is done by setting multiple
  pointers to a single expression, which is executed only once if necessary.
  This is  \emblue{\em call-by-need} evaluation.
\vspace*{-1em}
\end{itemize}
\item Eager evaluation can lead to \emred{non-terminating} computations for programs that necessarily terminate with 
  lazy evaluation.
\vspace*{-0.5em}
\[
(\lambda x : \tnat \ldot 0) ((\cfix \lambda y : \tnat \ldot y) 0)
\lazyreduce 0
\vspace*{-1em}
\]
\begin{eqnarray*}
& & (\lambda x : \tnat \ldot 0) ((\cfix \lambda y : \tnat \ldot y) 0)
\\ & \eagerreduce &
(\lambda x : \tnat \ldot 0) (((\lambda y: \tnat \ldot y) \delay{\cfix \lambda y: \tnat \ldot y} 0)
\\ & \eagerreduce &
(\lambda x : \tnat \ldot 0) ((\lambda z: \tnat \ldot (\cfix \lambda y: \tnat \ldot y) z) 0)
\\ & \eagerreduce &
(\lambda x : \tnat \ldot 0) ((\cfix \lambda y: \tnat \ldot y) 0)
\end{eqnarray*}
\end{itemize}
\end{slide}

\begin{slide}{Semantics}
\begin{itemize}
\item Roughly speaking, semantics is the meaning of programs. 
\\[0.4em]
\item Three styles of semantics
\begin{itemize}
\item \emblue{Operational semantics}: how do programs execute (reduce)? \\
  It's about the evaluation or computation procedure of programs.
\item \emblue{Axiomatic semantics}: do programs satisfy a certain property? \\
  It's about the reasoning of properties of and relations between programs.
\item \emblue{Denotational semantics}: what do programs represent? \\
  It's about the exact (mathematical) object that a program denotes.
\end{itemize}
\item In this course we only discuss operational semantics. For the other two semantics, 
  please refer to Prof. Liu Xinxin's course on formal semantics of programming languages.
\end{itemize}
\end{slide}

\begin{slide}{Operational semantics}
\begin{itemize}
\item The PCF reduction system indeed defines an operational semantics of PCF, which we often refer to as 
  \emblue{small-step} semantics.
\vspace*{-0.5em}
\item Big-step semantics: $e \Reduceto e'$ if $e \reduce^* e'$ and $\neg (e' \reduce)$.
\end{itemize}
\[\small
\begin{array}{c}
\begin{prooftree}
c \in \set{\ctrue, \cfalse, 0, 1, 2, \ldots}
\justifies
c \Reduceto c
\end{prooftree}
\qquad 
\begin{prooftree}
e_1 \Reduceto n_1 
\quad e_2 \Reduceto n_2 
\quad n_1 \;\mathtt{op}\; n_2  = n
% \quad \mathtt{op} \in \set{+, -, *, /}
\justifies
e_1 \cop e_2 \Reduceto n
\end{prooftree}
\\[1.8em]
\begin{prooftree}
e_1 \Reduceto n_1 
\quad e_2 \Reduceto n_2 
\quad n_1 \neq n_2
\justifies
e_1 \eqtest e_2 \Reduceto \cfalse 
\end{prooftree}
\qquad 
\begin{prooftree}
e_1 \Reduceto n 
\quad e_2 \Reduceto n 
\justifies
e_1 \eqtest e_2 \Reduceto \ctrue
\end{prooftree}
\qquad
\begin{prooftree}
e \Reduceto \ctrue 
\quad e_1 \Reduceto v_1 
\justifies
\cif{e}{e_1}{e_2} \Reduceto v_1
\end{prooftree}
\\[1.8em]
\begin{prooftree}
e \Reduceto \cfalse
\quad e_2 \Reduceto v_2
\justifies
\cif{e}{e_1}{e_2} \Reduceto v_2
\end{prooftree}
\qquad
\begin{prooftree}
e_1 \Reduceto v_1 
\quad e_2 \Reduceto v_2 
\justifies
\cpair{e_1, e_2} \Reduceto \cpair{v_1, v_2}
\end{prooftree}
\qquad 
\begin{prooftree}
e \Reduceto \cpair{v_1, v_2}
\justifies
\cproj_i \, e \Reduceto v_i
\end{prooftree}
\\[1.8em]
\begin{prooftree}
e \Reduceto v 
\justifies
\cinj_i e \Reduceto \cinj_i v
\end{prooftree}
\qquad
\begin{prooftree}
e \Reduceto \cinj_1 v_1
\quad e_1 [v_1 / x_1] \Reduceto v
\justifies
\mathtt{case} \; e \; \mathtt{of} \; \cinj_1 x_1 \; \mathtt{in} \; e_1 ; \cinj_2 x_2 \; \mathtt{in} \; e_2 
\Reduceto v
\end{prooftree}
\\[1.8em]
\begin{prooftree}
e \Reduceto \cinj_2 v_2
\quad e_2 [v_2 / x_2] \Reduceto v
\justifies
\mathtt{case} \; e \; \mathtt{of} \; \cinj_1 x_1 \; \mathtt{in} \; e_1 ; \cinj_2 x_2 \; \mathtt{in} \; e_2 
\Reduceto v
\end{prooftree}
\qquad
\begin{prooftree}
e_1 \Reduceto \lambda x \ldot e_1'
\quad e_1'[ e_2 / x ] \Reduceto v
\justifies
e_1 e_2 \Reduceto v
\using{\mbox{\emred{(Lazy)}}}
\end{prooftree}
%% \qquad
%% \begin{prooftree}
%% e (\cfix \, \delay{e}) \Reduceto v
%% \justifies
%% \cfix \, e \Reduceto v
%% \end{prooftree}
\end{array}
\]
\end{slide}

\begin{slide}{Total recursive functions}
\begin{itemize}
\item A class $\mathcal C$ of {\em numeric functions} $f: \vnat^k \to \vnat$ is
\vspace*{-0.5em} 
\small
\begin{itemize}
\item {\em closed under composition}: for every $f_1, \ldots , f_l: \vnat^k \to \vnat$ and 
$g: \vnat^l \to \vnat$ of $\mathcal C$, the function 
\empurple{$h(n_1, \ldots, n_k) = g(f_1(n_1, \ldots, n_k), \ldots, f_l(n_1, \ldots, n_k))$} is 
also in $\mathcal C$.
\item {\em closed under primitive recursion}: for every $f: \vnat^k \to \vnat$ and 
  $g: \vnat^{k+2} \to \vnat$ of $\mathcal C$, the function $h: \vnat^{k+1} \to \vnat $ 
  defined as follows is in $\mathcal C$: \vspace*{-0.5em}
\empurple{
\begin{eqnarray*}
h (0, n_1, \ldots , n_k) & = & f (n_1, \ldots, n_k) \\
h (m+1, n_1, \ldots, n_k) & = & g (m, n_1, \ldots, n_k, h(m, n_1, \ldots, n_k)) \\[-3.5em]
\end{eqnarray*}
}
\item {\em closed under minimization}: for every $f: \vnat^{k+1} \to \vnat$ of $\mathcal C$ 
  such that $\forall n_1, \ldots, n_k, \exists m, f(m, n_1, \ldots, n_k) = 0$, 
  the function \empurple{$g (n_1, \ldots, n_k) = \min(\set{m \mid f(m, n_1, \ldots, n_k) = 0})$}
  is in $\mathcal C$.
\vspace*{-0.5em}
\end{itemize}
\normalsize
\item \emblue{Total recursive functions} are the least class of numeric functions:
\vspace*{-0.5em}
\begin{itemize}
\item closed under composition, primitive recursion and minimization;
\item contains projection function \empurple{$\pmb{proj}_i^k (n_1, \ldots, n_k) = n_i$}, 
  successor function \empurple{$\pmb{succ}(n) = n+1$} 
  and constant zero function \empurple{$\pmb{zero}(n) = 0$}.
\end{itemize}
\end{itemize}
\end{slide}

\begin{slide}{Expressing total recursive functions in PCF}
\begin{itemize}
\item A numeric function $f: \vnat^k \to \vnat$ is {\em PCF-definable} 
  if there exists a PCF program $e$ of $\tnat^k \to \tnat$ such that 
  $\Eval{e} = f$ in the set-theoretic model.
\item Every total recursive function is definable in PCF. \\
{\em Proof.} It suffices to define a PCF term for every basic function and composition, primitive recursion, 
minimization.
\vspace*{-0.5em}
\begin{itemize}
\item Basic functions: $\overline{\pmb{proj}_i^k} \defeq \lambda x_1 \ldot \cdots \lambda x_k \ldot x_i$.
\item Composition: 
\vspace*{-0.5em}
\[
\overline{h} \defeq \lambda x_1 \ldot \cdots \lambda x_k \ldot 
\overline{g} (\overline{f_1} x_1 \cdots x_k) \cdots (\overline{f_l} x_1 \cdots x_k) 
\vspace*{-2.5em}
\]
\item Primitive recursion
\vspace*{-0.5em}
\[
\begin{array}{ll}
\overline{h} \defeq \! \! & \cfix \, \lambda h' \ldot \lambda n \ldot 
\\ &
\qquad\cif{n > 0}
{\lambda \vec{x} \ldot \overline{g} (n-1, \vec{x}, h' (n-1, \vec{x}))}
{\lambda \vec{x} \ldot \overline{f} \vec{x}}
\end{array}
\vspace*{-1.5em}
\]
\item Minimization: 
\vspace*{-0.5em}
\[
\overline{g} \defeq (\cfix \, \lambda g' \ldot \lambda n \ldot \lambda \vec{x} \ldot  
\cif{f (n, \vec x) \eqtest 0}{n}{g' (n+1, \vec x)}) 0
\vspace*{-0.5em}
\]
\end{itemize}
\end{itemize}
\end{slide}

\begin{slide}{Partial recursive functions}
Partial numeric functions $f: \vnat^k \to \vnat$ can be undefined at some points.
\vspace*{-0.5em}
\small
\begin{itemize}
\item {\em Closed under composition}: for every $f_1, \ldots , f_l: \vnat^k \to \vnat$ and 
$g: \vnat^l \to \vnat$ of $\mathcal C$, the composition function is in $\mathcal C$:
\vspace*{-0.6em}
\empurple{
\[
h(n_1, \ldots, n_k) = \left \{ 
\begin{array}{ll}
g(m_1, \ldots, m_l) & \mbox{if $m_i = f_i(n_1, \ldots, n_k)$ and}
\\ & \mbox{$g(m_1, \ldots, m_l)$ are defined}
\\[0.2em]
\bot & \mbox{otherwise}
\end{array}
\right.
\vspace*{-1.2em}
\]
}
\item {\em closed under primitive recursion}: for every $f_1, \ldots , f_l: \vnat^k \to \vnat$ and 
  $g: \vnat^{k+2} \to \vnat$ of $\mathcal C$, the function $h: \vnat^{k+1} \to \vnat $ 
  defined as follows is in $\mathcal C$: \vspace*{-0.5em}
\empurple{
\begin{eqnarray*}
h (0, n_1, \ldots , n_k) & = & 
\left \{ 
\begin{array}{ll}
f (n_1, \ldots, n_k) & \mbox{if $f(n_1, \ldots, n_k)$ is defined} \\
\bot & \mbox{o.w.}
\end{array}
\right .
\\
h (m+1, n_1, \ldots, n_k) & = & \left \{
\begin{array}{ll}
g (m, n_1, \ldots, n_k, q) & \mbox{if $q = h(m, n_1, \ldots, n_k)$ and}
\\ & \mbox{$g(m, n_1, \ldots , n_k, q)$ are defined} 
\\
\bot & \mbox{o.w.}
\end{array}
\right .
\\[-2.5em]
\end{eqnarray*}
}
\item {\em closed under minimization}: for every $f: \vnat^{k+1} \to \vnat$ of $\mathcal C$, 
  $\mathcal C$ contains the function 
\vspace*{-0.5em}
\empurple{
\[
g (n_1, \ldots, n_k) = \left \{
\begin{array}{l}
\mbox{the least $m$ s.t. $f(m, n_1, \ldots, n_k) = 0$ if such $m$ exist}  \\[0.2em]
\bot \mbox{ otherwise}
\end{array}
\right.
\]
}
\end{itemize}
\end{slide}

\begin{slide}{Expressing partial recursive functions in PCF}
\begin{itemize}
\item \emblue{Partial recursive functions} are the least class of partial numeric functions:
\vspace*{-0.5em}
\begin{itemize}
\item closed under composition, primitive recursion and minimization;
\item contains projection function \empurple{$\pmb{proj}_i^k$}, successor function \empurple{$\pmb{succ}$}
  and constant zero function \empurple{$\pmb{zero}$}.
\end{itemize}
\item Every partial recursive function is definable in PCF. \\[0.3em]
{\em Proof. } Representing undefinedness by \emred{non-termination} in PCF: 
\[
\overline{\bot} \defeq \cfix \, \lambda f \ldot f
\vspace*{-1.5em}
\]
\item By Church's thesis, all mechanically computable functions on natural numbers are definable in PCF.
\end{itemize}
\end{slide}

\begin{slide}{Non-definability of parallel operations}
\begin{itemize}
\item The {\em parallel-or} operation:
\[
\pmb{POR} \, e_1 \, e_2 \reduceto \left \{
\begin{array}{ll}
\ctrue & \mbox{if $e_1 \reduceto \ctrue$ \emred{or}  $e_2 \reduceto \ctrue$} 
\\
\cfalse & \mbox{if $e_1 \reduceto \cfalse$ and $e_2 \reduceto \cfalse$} 
\\
\mbox{no normal form} & \mbox{otherwise}
\end{array}
\right .
\]
where $e_1, e_2$ are closed PCF terms. 
\item There is no PCF program that defines $\pmb{POR}$.
\\[0.3em]
The proof can be found in Chapter 2 of J. Mitchell's book {\em Foundations for Programming Languages},
which relies essentially on the \emred{sequentiality} of PCF reduction system.
%% no matter what evaluation strategy it is, given an evaluation context $E\termhole$, 
%% the evaluation of the program $E[\pmb{POR} \, e_1 \, e_2]$ 
%% \vspace*{-0.5em}
%% \begin{itemize}
%% \item either does not evaluate $\pmb{POR} \, e_1 \, e_2$ at all, 
%% \item or eventually reduces to $E'[e_i]$ for $i = 1 , 2$ and another evaluation context $E'\termhole$. 
%% \end{itemize}
%% This says that the evaluation will never computes $e_1$ and $e_2$ at the same time.
\end{itemize}
\end{slide}

\end{document}
 

